梯度检查通过了，训练依然失败——这是本单元要处理的全部情形。反向传播只给出当前参数处的梯度，它不告诉你走多远仍然安全，不保证每个中间量都落在浮点数能表示的窗口里，也不承诺 mini-batch 采样出的那个带噪估计在单步上是下降方向。深度复合还让信号与梯度跨层按指数放大或衰减。于是学习率、初始化、优化器、归一化与随机丢弃这些看起来零散的旋钮，实际共同定义了一个动态系统。

本单元讲的是训练现场，不是收敛定理。凸优化那一部分的非凸单元负责数学侧——下降引理在非凸地形中的退化、鞍点的二阶判别、PL 条件、随机梯度留下的噪声平台；这里负责的是另一半：梯度爆炸与消失怎样在日志里看出来，warmup 在缓解什么，学习率调度背后的经验规律，混合精度下损失缩放挡住了哪一段浮点范围。两条线在同一批现象上交汇，问的问题不同。

三篇文章各自立一个判断。第一篇把训练拆成数学目标、随机采样与有限精度算术三方，说明深度让尺度按指数漂移而浮点窗口有限，并给出一条从小样本过拟合开始的诊断顺序。第二篇拆开动量与 Adam 各自在平均什么：前者平均梯度方向，后者还平均梯度平方，而梯度平方度量的是波动幅度而非曲率，因此逐坐标除法买到的只是对角子群上的尺度不变性，是 Newton 尺度不变性的一个廉价模仿。第三篇转向被优化的那张曲面：归一化改写损失曲面的条件数与梯度对权重尺度的响应，Dropout 在训练时定义了一族子网络，两者都让训练态与推理态成为两个不同的函数，而由此产生的故障通常不报错。

贯穿三篇的判断是同一条：深度模型把选择表示变成了可优化的对象，换来这一点的是几乎所有保证都退化为经验规律。这些机制改变的是时间平均、坐标缩放、有效步长与训练时的随机性，它们不能替代学习率检查、数据划分与独立验证。流传的训练配方只是这些机制的一种组合，换个任务未必仍然成立。

\subsection{怎么读这一章}

第一篇是诊断主干，训练出现 NaN、停滞或震荡时从这里进入。第二篇在需要比较 SGD、动量与 Adam，或者需要解释权重衰减与 \(L_2\) 惩罚为何在自适应方法里分家时回补。第三篇面向含归一化或 Dropout 的网络，部署之前必须确认训练态与推理态的差异有多大。

\subsection{本章篇目}

以下篇目按概念依赖排列；若从具体问题进入，也可以先打开最接近当前困惑的一篇，再沿篇内链接回补。

\begin{itemize}
\tightlist
\item \hyperref[sec:14-Deep-Learning/03-Optimization-and-Training/01]{``训练动力学与数值稳定''}；
\item \hyperref[sec:14-Deep-Learning/03-Optimization-and-Training/02]{``动量与 Adam 在平均什么''}；
\item \hyperref[sec:14-Deep-Learning/03-Optimization-and-Training/03]{``归一化与 Dropout 改变训练时的随机系统''}。
\end{itemize}

三篇读完，你手上会有一套能定位故障的监控量与一组关于旋钮的判断，但不会有关于``这个配置一定收敛''的承诺——那种承诺在这个层面上并不存在。
